\section{Related Work}
\label{sec:related}

\textit{Proof Logging for Symbolic Execution.}
%The proof-logging approach proposed by Trabish et al.~\cite{klee-rocq} is restricted to a subset of \llvm with integers.
Trabish et al.~\cite{klee-rocq} propose a proof-logging approach based on symbolic execution for a subset of \llvm with integers.
Lundberg et al.~\cite{lundberg2024proof} use symbolic execution to produce functional correctness proofs for P4 (a DSL for network devices) programs using HOL4.
% TODO: no memory modeling...
%They devise a symbolic execution algorithm that outputs a collection of theorems that can be used to prove properties about the analysis results.
In contrast, our approach extends support to memory operations, thus making it applicable to more realistic programs.

\textit{Verified \SE Engines.}
% TODO: shorten
Correnson et al.~\cite{fse23wise} develop a formally verified symbolic bug finder for a toy WHILE-style language in \coq.
Their formalization does not support memory operations,
and it is coupled with a specific search heuristic (BFS).
L\"{o}\"{o}w et al.~\cite{cse-oopsla25} mechanized a memory-parametric framework in \coq for a \textit{while}-like IR language,
and instantiated it on several memory models.
Based on this framework,
they built a prototype \SE engine that provides soundness and completeness guarantees.
However,
their \SE engine lacks support for features like recursion and has not been evaluated on real-world programming languages.
Adapting their core theory to \llvm IR requires modifying the underlying IR language (and core theory) and formalizing the memory model of \llvm IR,
which would demand a significant effort.
There exist \SE engines~\cite{gillian} (for C and Java) which are based on a similar theory,
but they are unverified and do not produce machine-checked proofs.
Keuchel et al.~\cite{icfp22:keuchel} built a mechanized \SE engine for a simplified subset of SAIL~\cite{sail},
but it lacks support for core memory features (\eg pointers and dynamic allocations) required by real-world programming languages.
%Katamaran~\cite{icfp22:keuchel} is a mechanized \SE engine for a simplified subset of SAIL~\cite{sail},
%but it lacks support for core memory features (\eg pointers and dynamic allocations) required by real-world programming languages.

%\textit{Auto-Generated \SE Engines.}
%Gillian~\cite{gillian} is a memory-parametric framework for constructing \SE engines that has been instantiated to JavaScript and C.
%However, the resulting \SE engines are unverified as they are implemented in OCaml,
%leaving a trust gap with respect to the intended semantics.
%L\"{o}\"{o}w et al.~\cite{cse-oopsla25} extend the framework in~\cite{gillian} and mechanize it in \coq,
%while instantiating it on a broad collection of memory models.
%However, they do not provide an executable implementation at all, \ie an \SE engine,
%as their work is focused on theoretical foundations.

\textit{Deductive Verification.}
% TODO: needed?
%The approach proposed by~\cite{icfp22:keuchel} is similar to the work mentioned above,
%but based on axiomatic semantics instead.
Veil~\cite{veil} is a foundational program verifier developed in Lean~\cite{moura2021lean}.
It provides automated verification for decidable fragments of first-order logic,
falling back to an interactive proof mode when automation fails.
It was evaluated on concurrent and distributed algorithms.
% LOOM
Loom~\cite{loom-popl26} is a framework for constructing foundational verifiers.
It integrates SMT-based proof automation with an interactive fallback mechanism.
The derived program verifier provides support for proof automation using SMT solvers, similarly to~\cite{veil}.
Loom was evaluated via Velvet, a Dafny-style verifier for a prototype imperative language with support for arrays.
However, Loom has yet to be demonstrated on real-world programming languages with full memory support (\eg \llvm IR).
To the best of our knowledge,
extending Loom to support such programming languages requires non-trivial changes to its underlying theoretical foundations.

\textit{Proof Logging for SMT.}
Modern SMT solvers~\cite{armand2011modular,ekici2017smtcoq,de2008proofs,marx2013proof,flippo2024multi,mcilree2024proof,van2025efficient} increasingly rely on proof logging to guarantee correctness.
Rather than blindly trusting the SMT solver, a proof-logging framework forces it to produce a formal justification of its reasoning.
Proof logging for SMT solvers is orthogonal to our approach and can be combined with it.
